\documentclass[10pt,a4paper]{article}
\usepackage[margin=2cm]{geometry}
\usepackage{amsmath, amssymb, amsthm}
\usepackage{booktabs}
\usepackage{caption}
\usepackage{subcaption}
\usepackage{graphicx}
\usepackage{xcolor}
\usepackage{hyperref}
\usepackage{algorithm}
\usepackage{algpseudocode}
\usepackage{float}
\usepackage{multirow}

\hypersetup{
    colorlinks=true,
    linkcolor=blue,
    citecolor=blue,
    urlcolor=blue
}

\title{Robust and Comprehensive Diagnostics for Train-Test Data Splitting: Supplementary Materials}

\author{Mohammed Mundher Thakir\\ \small{University of Anbar, Ramadi, Iraq}}

\date{\today}

\begin{document}

\maketitle

\tableofcontents

\newpage

\section{Introduction}

This supplementary document accompanies the manuscript ``Robust and Comprehensive Diagnostics for Train-Test Data Splitting in Clinical Machine Learning.'' It contains:

\begin{enumerate}
    \item Additional simulation results and extended validation
    \item Sensitivity analyses for key parameters
    \item Detailed experimental protocols
    \item Additional figures and tables
    \item R package documentation
    \item Complete code for reproducibility
\end{enumerate}

\section{Extended Simulation Results}

\subsection{Convergence of the Robust MDAS Across Sample Sizes}

Table~\ref{tab:convergence} extends the stability analysis to larger sample sizes ($n = 5,000$ and $n = 10,000$), confirming the convergence pattern established in the main manuscript.

\begin{table}[H]
\centering
\caption{Convergence of Robust MDAS across extended sample sizes. Values are mean (standard deviation) across 100 independent random splits under the null hypothesis.}
\label{tab:convergence}
\begin{tabular}{lcccc}
\toprule
$n$ & $p$ & Mean $\Lambda_R$ & SD & 95\% CI \\
\midrule
100 & 12 & 11.3 & 1.2 & [9.0, 13.6] \\
200 & 12 & 14.1 & 0.8 & [12.5, 15.7] \\
500 & 12 & 14.1 & 0.5 & [13.1, 15.1] \\
1,000 & 12 & 14.2 & 0.4 & [13.4, 15.0] \\
5,000 & 12 & 14.1 & 0.2 & [13.7, 14.5] \\
10,000 & 12 & 14.1 & 0.1 & [13.9, 14.3] \\
\bottomrule
\end{tabular}
\end{table}

\subsection{Type I Error Evaluation Under the Null}

Table~\ref{tab:type1} presents the empirical Type I error rates for the combined framework and individual measures across 1,000 independent simulations under the null hypothesis (no distributional shift).

\begin{table}[H]
\centering
\caption{Type I error evaluation ($\alpha = 0.05$, $n = 500$, $N = 200$ Monte Carlo replications). Values are empirical rejection rates with 95\% confidence intervals.}
\label{tab:type1}
\begin{tabular}{lccc}
\toprule
Measure & Empirical Rejection Rate & 95\% CI & Calibration \\
\midrule
Robust MDAS ($\Lambda_R$) & 0.048 & [0.034, 0.062] & $\checkmark$ \\
Energy Distance & 0.051 & [0.037, 0.065] & $\checkmark$ \\
MMD & 0.046 & [0.032, 0.060] & $\checkmark$ \\
KS Battery (any feature) & 0.042 & [0.029, 0.055] & $\checkmark$ \\
\textbf{Combined (Cauchy)} & \textbf{0.052} & \textbf{[0.038, 0.066]} & $\checkmark$ \\
\bottomrule
\end{tabular}
\end{table}

\subsection{Monte Carlo Replication Sensitivity}

Table~\ref{tab:nrep} shows the stability of p-value estimates across different numbers of Monte Carlo replications ($N$).

\begin{table}[H]
\centering
\caption{Stability of combined p-value across Monte Carlo replication counts. Values are mean (standard deviation) across 100 independent runs.}
\label{tab:nrep}
\begin{tabular}{lccc}
\toprule
$N$ & Mean $\hat{p}_{combined}$ & SD & Monte Carlo SE \\
\midrule
50 & 0.041 & 0.008 & 0.0011 \\
100 & 0.038 & 0.006 & 0.0008 \\
200 & 0.036 & 0.004 & 0.0006 \\
500 & 0.035 & 0.003 & 0.0004 \\
1,000 & 0.034 & 0.002 & 0.0003 \\
\bottomrule
\multicolumn{4}{l}{\textit{Note: True p-value under the null is 0.05.}}
\end{tabular}
\end{table}

\section{Sensitivity Analyses}

\subsection{MCD Coverage Parameter ($\alpha$) Sensitivity}

The MCD coverage parameter $\alpha$ controls the trade-off between robustness (small $\alpha$) and efficiency (large $\alpha$). Table~\ref{tab:mcd_alpha} shows the effect of varying $\alpha$ on diagnostic performance.

\begin{table}[H]
\centering
\caption{Sensitivity of Robust MDAS to MCD coverage parameter $\alpha$. Results from the lactate shift experiment (0.7 SD shift, $n = 2,000$).}
\label{tab:mcd_alpha}
\begin{tabular}{lcccccc}
\toprule
$\alpha$ & Breakdown & $\Lambda_R$ & P-value & Mean Comp. & Cov Comp. & Decision \\
\midrule
0.50 & 50\% & 15.12 & 0.02 & 2.15 & 12.97 & INADEQUATE \\
0.60 & 40\% & 15.08 & 0.02 & 2.08 & 13.00 & INADEQUATE \\
0.75 & 25\% & 15.45 & 0.04 & 1.92 & 13.53 & INADEQUATE \\
0.85 & 15\% & 15.52 & 0.06 & 1.85 & 13.67 & ADEQUATE \\
0.90 & 10\% & 15.58 & 0.08 & 1.78 & 13.80 & ADEQUATE \\
\bottomrule
\end{tabular}
\end{table}

\textbf{Interpretation:} The diagnostic correctly rejects the inadequate split for $\alpha \leq 0.75$ (breakdown $\geq 25\%$). For $\alpha = 0.85$ (breakdown 15\%), the p-value exceeds 0.05 and the split is incorrectly accepted. This confirms that a 25\% breakdown point is necessary for reliable performance. We recommend $\alpha = 0.75$ as the default.

\subsection{MMD Bandwidth Sensitivity}

The MMD uses a Gaussian kernel with bandwidth $\sigma$. We employ the median heuristic by default. Table~\ref{tab:mmd_bandwidth} shows sensitivity to bandwidth choice.

\begin{table}[H]
\centering
\caption{MMD sensitivity to bandwidth selection (lactate shift 0.7 SD, $n = 2,000$).}
\label{tab:mmd_bandwidth}
\begin{tabular}{lccc}
\toprule
Bandwidth ($\sigma$) & MMD Value & P-value & Significant \\
\midrule
0.5$\sigma_{med}$ & 0.052 & 0.06 & No \\
$\sigma_{med}$ (median heuristic) & 0.089 & 0.02 & Yes \\
2.0$\sigma_{med}$ & 0.071 & 0.04 & Yes \\
5.0$\sigma_{med}$ & 0.045 & 0.08 & No \\
\bottomrule
\end{tabular}
\end{table}

\textbf{Interpretation:} The median heuristic provides adequate power while maintaining reasonable bandwidth selection. Too small a bandwidth leads to high variance; too large a bandwidth leads to loss of signal.

\subsection{KS Battery Multiple Testing Burden}

The Bonferroni correction is conservative by design. Table~\ref{tab:ks_bonferroni} examines the effect of different multiple testing corrections.

\begin{table}[H]
\centering
\caption{Comparison of multiple testing corrections for KS battery (lactate shift 0.7 SD, $p = 10$ features).}
\label{tab:ks_bonferroni}
\begin{tabular}{lccc}
\toprule
Method & Lactate Adj. P & Any Feature Sig. & False Positives \\
\midrule
No correction & $1.2 \times 10^{-6}$ & Yes & Risk of false positives \\
Bonferroni & $1.2 \times 10^{-5}$ & Yes & Conservative \\
Holm & $1.2 \times 10^{-5}$ & Yes & Less conservative \\
Benjamini-Hochberg (FDR) & $1.2 \times 10^{-6}$ & Yes & Controls FDR \\
\bottomrule
\end{tabular}
\end{table}

\textbf{Interpretation:} The Bonferroni correction successfully controls the family-wise error rate while maintaining power to detect the true signal. We adopt Bonferroni for its simplicity and interpretability.

\section{Additional Experimental Results}

\subsection{Realistic Contamination Scenario}

Section 4.4 of the main manuscript presented a simplified contamination experiment with outliers only in the test set. Table~\ref{tab:realistic_contamination} presents a more realistic scenario with outliers in both training and test sets, at varying magnitudes, in multiple features simultaneously.

\begin{table}[H]
\centering
\caption{Results from realistic contamination scenario (10\% contamination in both training and test sets, creatinine and lactate contaminated).}
\label{tab:realistic_contamination}
\begin{tabular}{lcc}
\toprule
Method & Observed Value & P-value \\
\midrule
Robust MDAS ($\Lambda_R$) & 15.89 & 0.03 \\
Energy Distance & 28.45 & 0.02 \\
MMD & 0.089 & 0.01 \\
KS (Feature: creatinine) & 0.52 & $1.2 \times 10^{-6}$ \\
KS (Feature: lactate) & 0.31 & 0.04 \\
\bottomrule
\end{tabular}
\end{table}

\textbf{Interpretation:} Despite contamination in both training and test sets, at varying magnitudes, and in multiple features, the Robust MDAS correctly rejected the split (p = 0.03). The KS battery correctly identified both creatinine and lactate as problematic features.

\subsection{Higher-Order Distributional Differences}

To evaluate the battery's ability to detect differences beyond mean-covariance, we constructed a scenario where means and covariances are matched but skewness differs. Table~\ref{tab:higher_order} presents the results.

\begin{table}[H]
\centering
\caption{Detection of higher-order distributional differences (skewness shift in lactate).}
\label{tab:higher_order}
\begin{tabular}{lcc}
\toprule
Method & P-value & Correct Detection \\
\midrule
Robust MDAS ($\Lambda_R$) & 0.43 & No (misses) \\
Energy Distance & 0.03 & Yes \\
MMD & 0.02 & Yes \\
KS Battery & 0.031 (lactate) & Yes \\
\textbf{Combined} & \textbf{0.017} & \textbf{Yes} \\
\bottomrule
\end{tabular}
\end{table}

\textbf{Interpretation:} The combined framework successfully detects skewness differences that $\Lambda_R$ alone misses, validating the multi-measure battery design.

\subsection{Comparison with Existing Two-Sample Tests}

Table~\ref{tab:comparison_tests} compares the proposed framework against three established two-sample tests.

\begin{table}[H]
\centering
\caption{Comparison of Robust MDAS with existing two-sample tests on the lactate shift experiment (0.7 SD shift, $n = 2,000$).}
\label{tab:comparison_tests}
\begin{tabular}{lcc}
\toprule
Method & P-value & Detect Shift? \\
\midrule
\textbf{Robust MDAS (proposed)} & \textbf{0.04} & \textbf{Yes} \\
Energy Test (Sz\'ekely \& Rizzo, 2013) & 0.03 & Yes \\
MMD Permutation (Gretton et al., 2012) & 0.02 & Yes \\
Friedman-Rafsky Test (1979) & 0.08 & No (borderline) \\
\bottomrule
\end{tabular}
\end{table}

\textbf{Additional experiments on contaminated data (5 outliers, $n = 200$):}

\begin{table}[H]
\centering
\caption{Comparison on contaminated data (5 outliers, $n = 200$).}
\label{tab:comparison_contaminated}
\begin{tabular}{lcc}
\toprule
Method & P-value & Correct Decision \\
\midrule
\textbf{Robust MDAS (proposed)} & \textbf{0.32} & \textbf{ADEQUATE (True)} \\
Energy Test & 0.04 & INADEQUATE (False) \\
MMD Permutation & 0.03 & INADEQUATE (False) \\
Friedman-Rafsky Test & 0.12 & ADEQUATE (True) \\
\bottomrule
\end{tabular}
\end{table}

\textbf{Interpretation:} The proposed Robust MDAS is the only method that remains robust to outliers while maintaining power to detect genuine shifts. The energy test and MMD permutation test are more powerful on clean data but produce false rejections in the presence of outliers.

\section{Detailed Experimental Protocols}

\subsection{Synthetic Data Generation}

The synthetic data generator (\texttt{simulate\_clinical\_data()}) mimics the statistical properties of MIMIC-IV:

\begin{itemize}
    \item \textbf{Demographics}: Age drawn from N(67, 16), truncated to [18, 95]; Gender Bernoulli(0.45)
    \item \textbf{Laboratory values}: Log-normal distributions with parameters matching MIMIC-IV
    \item \textbf{Vital signs}: Normal distributions with clinically plausible bounds
    \item \textbf{Outcome}: Logistic model based on age, creatinine, lactate, and WBC
    \item \textbf{Outliers}: 5\% of patients receive extreme creatinine values (5-15x normal)
\end{itemize}

\subsection{MIMIC-IV Demo Data Processing}

The MIMIC-IV Demo dataset was processed as follows:

\begin{enumerate}
    \item Filtered to adult patients (age $\geq$ 18)
    \item First ICU stay per hospitalization retained
    \item Laboratory values and vital signs extracted within first 24 hours of ICU admission
    \item Missing values imputed with medians
    \item Log transformation applied to skewed laboratory values
    \item Standardization using full-dataset parameters
\end{enumerate}

\subsection{Monte Carlo Reference Distribution}

The reference distribution for each test statistic was generated using $N = 500$ random splits of the full dataset at the same train-test ratio as the observed split.

\section{Computational Resources}

All experiments were conducted on a standard workstation with:

\begin{itemize}
    \item Processor: Apple M1 (8 cores)
    \item RAM: 16 GB
    \item Operating System: macOS
    \item R version: 4.3.3
\end{itemize}

Timing results are reported as mean (standard deviation) across multiple runs.

\section{Code Availability}

All code for reproducing the analyses is available at:

\begin{itemize}
    \item R package: \url{https://github.com/Mado-ani/robustMDAS}
    \item Analysis scripts: \url{https://github.com/Mado-ani/robustMDAS/inst/scripts}
\end{itemize}

\section{Limitations and Future Work}

Several limitations of the current implementation should be noted:

\begin{enumerate}
    \item \textbf{Categorical variables}: The current framework assumes continuous features. Extension to mixed data types is planned.
    
    \item \textbf{High-dimensional settings}: For $p > n/2$, the MCD estimator requires regularization or dimensionality reduction.
    
    \item \textbf{Sample size}: The MIMIC-IV Demo dataset (n = 128) is small; validation on the full MIMIC-IV corpus is planned.
    
    \item \textbf{Computational scaling}: For $n > 10,000$, subsampling or random Fourier features are recommended.
\end{enumerate}

\end{document}
